\documentclass[12pt]{article}
\usepackage[utf8]{inputenc}
\usepackage[T1]{fontenc}
\usepackage{hyperref}
\title{Mbangwe Language Resources for K-12 Curriculum Development}
\author{Compiled by Claude Sonnet 5.5 (claude-sonnet-5-5) for Andr\'as Kornai}
\date{}
\begin{document}
\maketitle

\section{Introduction}
Mbangwe (ISO 639-3 zmn, Glottocode mban1268) is a Bantu language of the Atlantic-Congo family. It is spoken in Gabon, mainly in Haut-Ogoou\'e Province near Franceville, and in Congo. Sources do not agree on which Congo is meant. Speaker estimates are about 3,500 to 4,300. Other names are M'bahouin and Mbaouin. Web search found only general reference and demographic pages. It found no school materials, primers, dictionaries or grammars.

\section{Primary Educational Resources}
No primers, textbooks or school programmes for Mbangwe turned up in search. The language is not known to be taught in schools.

\section{Linguistic and Reference Materials}
\begin{itemize}
\item Wikipedia, Mbangwe language: \url{https://en.wikipedia.org/wiki/Mbangwe_language}
\item Glottolog, Mbangwe (mban1268): \url{https://glottolog.org/resource/languoid/id/mban1268.bigmap.html}
\item ASJP, Mbangwe word list: \url{https://asjp.clld.org/languages/MBANGWE}
\end{itemize}

\section{Religious and Community Resources}
\begin{itemize}
\item Joshua Project, Mbangwe (Mbahouin) in Gabon: \url{https://www.joshuaproject.net/index.php/people_groups/13619/GB}
\item Joshua Project, Mbangwe in the Republic of the Congo: \url{https://www.joshuaproject.net/index.php/people_groups/13619/CF}
\item Joshua Project, language page for zmn: \url{https://www.joshuaproject.net/languages/zmn}
\item Scripture Earth page for zmn: \url{https://www.scriptureearth.org/00eng.php?iso=zmn}
\end{itemize}
I did not check what scripture the Scripture Earth page lists.

\section{Digital and Audio Resources}
Search found no audio, video or corpus resources for Mbangwe.

\section{Curriculum Development Context}
Documentation is thin, and there are no known teaching materials. Curriculum work would have to start with fieldwork and community consultation. Standardising the orthography and building a basic lexicon would come first. The ASJP word list is a small starting point for the lexicon.

\end{document}